\begin{claim}{D.34}{Proved}{Global signs and the cusp cycle of a pure involution}
Every nonidentity finite-order element of the full pure
three-cusp loop-outer group $P$ of D.32 is an involution
whose common multiplier is the global element
$-1\in\widehat\Z^\times$.  In particular this holds in
$GT_0$, and the multiplier-one kernel in either group is
torsion-free.

For its normalized lift
\[
 \alpha(x)=x^{-1},\qquad\alpha(y)=f^{-1}y^{-1}f,
                      \qquad f\in F',
                                                        \tag{D.34a}
\]
there are precisely three $F$-conjugacy classes of sections
of $F\rtimes_\alpha C_2\to C_2$.  Each of their fixed
subgroups has zero first cohomology with $\F_2$ coefficients.
The two section classes in each cusp normalizer are the
endpoints of one edge of a triangle.  Moreover there is
$h\in F'$ with
\[
                         f=h^{-1}\alpha(h).               \tag{D.34b}
\]
Thus $\alpha$ inverts both $x$ and $hyh^{-1}$.
The endomorphism $x\mapsto x$, $y\mapsto hyh^{-1}$ induces
an automorphism of every maximal pro-$p$ quotient.  Its
surjectivity on the full $F$, and conjugacy of the original
involution to complex conjugation in the pure loop group,
are not conclusions of this statement.
\end{claim}
\begin{proof}
We first prove a genus-zero incidence calculation for an
arbitrary number $n\geq3$ of cusps.  Let
\[
 F_n=\widehat{\langle c_1,\ldots,c_n\mid
                              c_1\cdots c_n=1\rangle}
\]
and let $\beta$ be an actual involution preserving each
peripheral conjugacy class, with common multiplier $\mu$
and $\mu_2=-1$.  Put $E_n=F_n\rtimes_\beta C_2$.
The equality $\beta^2=1$ gives $\mu^2=1$; hence each
$\mu_q$ is $1$ or $-1$.

For each cusp $L_i=\overline{\langle c_i\rangle}$,
A.7 gives an exact sequence
\[
 1\longrightarrow L_i\longrightarrow N_{E_n}(L_i)
                          \longrightarrow C_2\longrightarrow1.
                                                        \tag{D.34c}
\]
Choose a lift $a_i$ of the nonidentity element and write
$a_i c_i a_i^{-1}=c_i^\mu$ and $a_i^2=c_i^{v_i}$.
Then $(\mu-1)v_i=0$.  In every coordinate with
$\mu_q=-1$ this forces $(v_i)_q=0$, including $q=2$.
In a coordinate with $\mu_q=1$ the prime $q$ is odd,
so division by two is possible.  Thus
$(1+\mu)u_i=-v_i$ has a solution in $\widehat\Z$,
and $b_i=a_i c_i^{u_i}$ is an involution.

The set of its further order-two translates by $L_i$,
modulo conjugation by $L_i$, is exactly
\[
 \frac{\ker(1+\mu:\widehat\Z\to\widehat\Z)}
      {(1-\mu)\widehat\Z}\ \simeq\ \F_2.               \tag{D.34d}
\]
At an odd prime with $\mu_q=-1$, multiplication by two
is onto, and a coordinate with $\mu_q=1$ contributes
zero to the kernel.  The remaining coordinate is
$\Z_2/2\Z_2$.  Representatives for the two classes are
$b_i$ and $c_i^{e_2}b_i$.

Since $\beta$ is pure and the cusp classes span homology,
its action on $V=H_1(F_n,\F_2)$ is the identity.  Therefore
the actual mod-two abelianization of the semidirect product is
\[
                  H_1(E_n,\F_2)=V\oplus\F_2.             \tag{D.34e}
\]
The two involutions just exhibited have images differing by
the nonzero vector $[c_i]$ in $V$.  They cannot be conjugate
in $E_n$ or in $F_n$.  Thus the two local classes in
(D.34d) give two distinct global section classes.

Build a finite graph with one edge for each of the $n$
cusps and with these two global classes as its endpoints.
This graph uses some of the global section classes; we
now bound all of them.  The Brown calculation of D.28,
with free rank $n-1$ and prime two, gives
\[
 \sum_{[Q]}
   \bigl(1+\dim_{\F_2}H^1(C_{F_n}(Q),\F_2)\bigr)
       =\dim_{\F_2}H^k(E_n,\F_2)\leq n\qquad(k>1).
                                                        \tag{D.34f}
\]
Every order-two subgroup is a section, and its total-group
and kernel conjugacy classes agree, because $E_n=F_nQ$.
Hence there are at most $n$ vertices available for the graph.
An edge has no loop by (D.34e).

Assign to every vertex its image $(v,1)$ in (D.34e).
The sum of the edge vectors around a graph cycle is zero,
since its successive vertex differences cancel.  But the
$n$ vectors $[c_1],\ldots,[c_n]$ have just one nonzero
linear relation over $\F_2$, namely
$[c_1]+\cdots+[c_n]=0$.  Consequently any simple cycle
in the graph uses all $n$ edges.  The graph cannot be a
forest: a forest with $n$ edges needs at least $n+1$
vertices, more than (D.34f) permits.  Thus it is exactly
a simple $n$-cycle.  It has $n$ vertices, so there are
no additional section classes outside this graph, and
(D.34f) gives
\[
            H^1(C_{F_n}(Q),\F_2)=0
                   \quad\text{for every section }Q.       \tag{D.34g}
\]
This construction is invariant under any automorphism
of the extension preserving its kernel and peripheral
classes, with any permutation of those classes: such
an automorphism sends the normalizer and its two classes
in (D.34c)--(D.34d) to those of the image cusp.

Return to a nontrivial finite-order element of $P$.
D.32 gives its order two, its normalized actual lift
$\alpha$, and $\lambda_2=-1$.  We show that
$\lambda_q=-1$ for every odd prime $q$.  The only other
possibility is $\lambda_q=1$.  Suppose this occurs and
use the exact cyclic quotient
\[
 \chi:F=\widehat F(x,y)\longrightarrow C_q,
          \qquad\chi(x)=1,\quad\chi(y)=0,
          \qquad N=\ker\chi.                            \tag{D.34h}
\]
The common multiplier gives
$\chi\alpha=\lambda\chi=\chi$, since its reduction
modulo $q$ is one.  In particular $N$ is invariant.
The corresponding degree-$q$ covering of the pair of
pants has one boundary over $x$, one over
$z=(xy)^{-1}$, and $q$ over $y$, with representative
generators
\[
 x^q,\qquad z^q,\qquad x^i yx^{-i}\quad(0\leq i<q).
                                                        \tag{D.34i}
\]
The two first degrees are $q$ and all others are one.
Its Euler characteristic is $-q$ and it has $q+2$
boundaries, so its genus is zero.  Finite-index
completion identifies its full profinite group with
the specified subgroup $N$.

The restriction of $\alpha$ is pure on this entire
peripheral structure, with the same multiplier
$\lambda$.  The first two boundaries are the unique
lifts of their original boundaries.  For the other
$q$ boundaries use the actual normalized formula
\[
 \alpha(x^i yx^{-i})
        =x^{\lambda i}f^{-1}y^\lambda f x^{-\lambda i}.
                                                        \tag{D.34j}
\]
Here $\chi(f)=0$, because $f\in F'$, and
$\lambda i\equiv i\pmod q$.  The conjugator in
(D.34j) therefore belongs to the same $N$-coset
as $x^i$: explicitly
$k_i=x^{\lambda i}f^{-1}x^{-i}\in N$ satisfies
$\alpha(x^i yx^{-i})=k_i(x^i yx^{-i})^\lambda k_i^{-1}$.
This is the required conjugacy inside $N$.  The action on
$H_1(N,\Z_2)$ is consequently the scalar $-1$,
so the restriction remains a nontrivial involution.

Put $E=F\rtimes_\alpha C_2$ and
$E_N=N\rtimes_\alpha C_2$.  The quotient maps give
\[
                     E/N=C_q\times C_2.                 \tag{D.34k}
\]
In particular $E_N$ is normal in $E$.  Conjugation by
$x$ defines an automorphism of $E_N$ preserving its
map to $C_2$, and thus acts on its $N$-conjugacy
classes of sections.  Its $q$th power is conjugation
by $x^q\in N$, so this action has order dividing
$q$ on those classes and on the cusp incidence graph.
On the cusp edges (D.34i), it fixes the $x^q$ and
$z^q$ classes and cycles the $q$ classes over $y$.
The second fixed class is the unique boundary over
$z$, or follows directly from its cyclic coset
decomposition.  Thus its graph action has order
exactly $q$.

By the preceding calculation the graph is a cycle
with $q+2$ edges.  Its automorphism group has order
$2(q+2)$: an image of a vertex and a choice of which
neighbor is its successor determine the automorphism.
The odd prime $q$ does not divide $2(q+2)$.
This contradicts the displayed order-$q$ graph action.
Therefore $\lambda_q=-1$ for every odd $q$, proving
the global multiplier assertion.  This argument used
the whole pure group $P$, without an additional
$GT_0$ symmetry assumption.

At $n=3$ the incidence graph is the stated triangle.
Write $s$ for the chosen section in $E=F\rtimes C_2$.
The normalized group law gives $f\alpha(f)=1$,
so $fs$ is also an involution and it normalizes $y$
by inversion.  The three involutions
\[
                     s,\quad x^{e_2}s,\quad y^{e_2}fs
                                                        \tag{D.34l}
\]
have the distinct images $(0,0,1)$, $(1,0,1)$,
and $(0,1,1)$ in (D.34e).  They represent all three
classes.  Since $f\in F'$, the image of $fs$ is
$(0,0,1)$, so it is conjugate to $s$ by an element
of $F$.  Choose $h$ with $fs=h^{-1}sh$.  This is
exactly (D.34b).  In abelianization it says
$0=(-2)[h]$ in the torsion-free group
$\widehat\Z^2$, hence $h\in F'$.
Substitution gives
\[
 \alpha(hyh^{-1})
   =(hf)(f^{-1}y^{-1}f)(f^{-1}h^{-1})
   =hy^{-1}h^{-1}.
\]
The two displayed inverted generators have the same
abelian images as $x,y$.  On each maximal pro-$p$
quotient their images therefore generate the Frattini
quotient and hence the whole group.  The induced
endomorphism of that finitely generated free pro-$p$
group is surjective and thus an automorphism.
No full-profinite surjectivity is inferred from those
separate comparisons.

The program \texttt{code/part\_c/pure\_involution\_d34.py}
checks the normalizer orbit calculation, the unique
binary cusp relation, and the exact cyclic-cover
boundary and deck permutations for several odd
primes.  The exclusion of all odd primes is the
unrestricted divisibility argument above.
\end{proof}
